\begin{alist}
\item Η συνάρτηση $f$ ορίζεται για εκείνα μόνο τα $x$ για τα οποία ισχύει $10^x - 2 > 0 \Leftrightarrow 10^x > 2 \Leftrightarrow x > \log 2$. Άρα το πεδίο ορισμού της συνάρτησης $f$ είναι το $A = (\log 2, +\infty)$.

\item 
\begin{rlist}
\item Με $x \in (\log 2, +\infty)$, έχουμε:
\begin{align*}
\sqrt{\frac{10^x}{3}} &= \sqrt{10^x - 2} \Leftrightarrow \\
\left(\sqrt{\frac{10^x}{3}}\right)^2 &= \left(\sqrt{10^x - 2}\right)^2 \Leftrightarrow \\
\frac{10^x}{3} &= 10^x - 2 \Leftrightarrow \\
10^x &= 3\cdot 10^x - 6 \Leftrightarrow \\
6 &= 2\cdot 10^x \Leftrightarrow \\
10^x &= 3 \Leftrightarrow x = \log 3.
\end{align*}

\item Για να βρούμε τα κοινά σημεία των γραφικών παραστάσεων των συναρτήσεων $f$ και $g$, επιλύουμε την εξίσωση $g(x) = f(x)$ για τα κοινά $x$ στα οποία αυτές ορίζονται, δηλαδή για $x \in (\log 2, +\infty)$. \\
Είναι:
\[g(x) = f(x) \Leftrightarrow \log\sqrt{\frac{10^x}{3}} = \log\sqrt{10^x - 2} \Leftrightarrow \sqrt{\frac{10^x}{3}} = \sqrt{10^x - 2} \overset{(\beta\text{i})}{\iff} x = \log 3\]
και
\[y = f(\log 3) = \log\sqrt{10^{\log 3} - 2} = \log\sqrt{3 - 2} = \log 1 = 0.\]
Ώστε οι γραφικές παραστάσεις των συναρτήσεων $f$ και $g$ έχουν ένα κοινό σημείο, το $(\log 3, 0)$.
\end{rlist}
\end{alist}
